\section{Separation from Hankel images}
\label{sec:separation}

Retain the Hankel matrix $T$, the spaces $K$, $R=K^\perp\subset P$,
and $W=L\oplus R$ from the preceding section.  In particular,
$\dim W=d_*=20$ and $H=W\oplus K$ orthogonally.  Let $\mathcal C$
be the finite set of projective classes of nonzero symmetric operators
on $W$ supplied by Lemma~\ref{lem:inverse-limits}.
We will construct a vector $w$ and a subspace $U$ such that $w$ is
outside the closure of all images of $U$ obtained by applying one
Hankel matrix componentwise.  A quadratic form separating $w$ from
those images will then give the required positive semidefinite
quadratic matrix $Q$.

Fix an integer $N>10h^2$, where $h=\dim H=36$, and put
\begin{equation}
\label{eq:separation-dimensions}
 r=d_*N-1.
\end{equation}
We equip $H^N$ with the sum of the Euclidean inner products on its
factors.  For a subspace $F\subset H$, write $P_F$ for orthogonal
projection onto $F$; the notation $P_F^N$ denotes its componentwise
action on $H^N$.
The two subspaces in the third condition below will represent
directions in which eigenvalues diverge or have finite nonzero
limits.  Their different projection constraints arise by taking
limits of bounded Hankel images of unit tuples.

\begin{lemma}[Choice of the vector and subspace]
\label{lem:generic-wu}
There exist a unit vector $w=(w_1,\ldots,w_N)\in W^N$ and an
$r$-dimensional subspace $U\subset w^\perp\subset H^N$ with the
following properties.
\begin{enumerate}
\item The vectors $w_1,\ldots,w_N$ span $W$, and
\begin{equation}
\label{eq:inverse-pairings}
 \sum_{i=1}^N w_i^T Cw_i\ne0
 \qquad\text{whenever }[C]\in\mathcal C.
\end{equation}
\item The map
\begin{equation}
\label{eq:projection-U}
 P_W^N\big|_U:U\longrightarrow w^\perp\cap W^N
\end{equation}
is an isomorphism.
\item Suppose that $F_\infty,F_0\subset H$ are orthogonal subspaces,
$C_0$ is a symmetric operator on $F_0$, and
$\dim(F_\infty\oplus F_0)>d_*$.  No nonzero tuple
$u=(u_1,\ldots,u_N)\in U$ satisfies
\begin{equation}
\label{eq:forbidden-projections}
 P_{F_\infty}u_i=0,\qquad
 P_{F_0}u_i=C_0P_{F_0}w_i
 \qquad(1\le i\le N).
\end{equation}
\end{enumerate}
\end{lemma}

\begin{proof}
For a nonzero symmetric operator $C$ on $W$, the polynomial
$\sum_i w_i^T Cw_i$ on $W^N$ is nonzero: a vector with just one
nonzero component already detects a nonzero value of the quadratic
form of $C$.  Its nonvanishing therefore defines a dense open subset
of $W^N$.  The spanning condition is also dense and open, since
$N\ge\dim W$.  There are only finitely many classes in
$\mathcal C$, and changing a representative does not affect
nonvanishing.  We may consequently choose $w$ satisfying all these
conditions and then normalize it to unit length.

Fix this $w$ and consider the Grassmannian
\[
 \mathcal G=\Gr(r,w^\perp),\qquad
 \dim w^\perp=hN-1.
\]
Inside $w^\perp$ we have the orthogonal decomposition
\[
 w^\perp=(w^\perp\cap W^N)\oplus K^N.
\]
The $r$-planes complementary to $K^N$ form a nonempty open subset
of $\mathcal G$.  They are exactly the subspaces satisfying
\eqref{eq:projection-U}.

We show that the subspaces violating the third property have smaller
dimension than $\mathcal G$.  Fix dimensions for $F_\infty$ and
$F_0$, and set $j=\dim F_\infty+\dim F_0>d_*$.  The triples
$(F_\infty,F_0,C_0)$ form a semialgebraic parameter space of
dimension at most $4h^2$.  For example, orthogonal projections onto
the two spaces and the extension of $C_0$ by zero give matrix
coordinates with this bound.  For each fixed triple,
\eqref{eq:forbidden-projections} prescribes the component of every
$u_i$ in $F_\infty\oplus F_0$.  Its solution set in $H^N$ is
therefore an affine space of dimension $(h-j)N$.

For a fixed nonzero $u\in w^\perp$, the subspaces in
$\mathcal G$ containing $u$ form a Grassmannian of codimension
\[
 (hN-1)-r=(h-d_*)N
\]
in $\mathcal G$.  Apply the semialgebraic dimension bounds to the
incidence relation consisting of a parameter triple, a nonzero
solution $u\in w^\perp$ of
\eqref{eq:forbidden-projections}, and an $r$-plane containing $u$.
Its image in $\mathcal G$ has dimension at most
\begin{align*}
 \dim\mathcal G+4h^2+(h-j)N-(h-d_*)N
 &=\dim\mathcal G+4h^2-(j-d_*)N\\
 &<\dim\mathcal G.
\end{align*}
There are finitely many possible pairs of dimensions for
$F_\infty,F_0$.  The union of their exceptional images still has
smaller dimension and cannot contain the nonempty open set of
complements to $K^N$.  Choosing $U$ outside that union proves both
remaining properties.
\end{proof}

Fix $w$ and $U$ as in Lemma~\ref{lem:generic-wu}.  Notice that the
$P^N$ component of $w$ is nonzero: otherwise all its components would
lie in $L$ and could not span $W$.  We now use the third property of
the lemma to control every possible limit of Hankel images, including
limits produced by unbounded matrices.

\begin{lemma}[Separation from the closure]
\label{lem:separation}
For $w$ and $U$ as in Lemma~\ref{lem:generic-wu}, let
\[
 \mathcal S_U=
 \{(I_N\otimes Y)u: u\in U,\ Y\text{ a Hankel matrix on }H\}.
\]
Then $w\notin\overline{\mathcal S_U}$.
\end{lemma}

\begin{proof}
Suppose, to the contrary, that there are tuples $u^{(n)}\in U$ and
Hankel matrices $Y_n$ such that
\[
 (I_N\otimes Y_n)u^{(n)}\longrightarrow w.
\]
After omitting finitely many terms the tuples are nonzero.  Replacing
$u^{(n)}$ by $u^{(n)}/\|u^{(n)}\|$ and $Y_n$ by
$\|u^{(n)}\|Y_n$ leaves the products unchanged.  We may thus assume
that $\|u^{(n)}\|=1$.  Pass to a subsequence for which
$u^{(n)}\to u\in U$, where $\|u\|=1$.

Choose an orthonormal eigenbasis $e_{1,n},\ldots,e_{h,n}$ for each
$Y_n$, with eigenvalues $\lambda_{1,n},\ldots,\lambda_{h,n}$.
After a further subsequence, each eigenvector converges and each
eigenvalue converges in the extended real line.  The limiting
eigenvectors form an orthonormal basis of $H$.  Let $F_\infty$ be
the span of those corresponding to eigenvalues whose absolute
values tend to infinity, and let $F_0$ be the span of those
corresponding to finite nonzero limits.  All remaining eigenvalues
tend to zero.  Define $C_0$ on $F_0$ by taking the reciprocal of
each corresponding eigenvalue limit in this limiting basis.

For each component of the tuples, the spectral identity is
\begin{equation}
\label{eq:spectral-components}
 \langle e_{\ell,n},Y_nu_i^{(n)}\rangle
 =\lambda_{\ell,n}\langle e_{\ell,n},u_i^{(n)}\rangle.
\end{equation}
When $\lambda_{\ell,n}\to0$, the right side tends to zero because
$u^{(n)}$ is bounded.  It follows that each $w_i$ belongs to
$F:=F_\infty\oplus F_0$, and the spanning property of $w$ gives
$W\subset F$.  When $|\lambda_{\ell,n}|\to\infty$, divide
\eqref{eq:spectral-components} by $\lambda_{\ell,n}$ to see that
the corresponding component of $u_i$ is zero.  For finite nonzero
limits, the same identity gives
\[
 P_{F_\infty}u_i=0,\qquad
 P_{F_0}u_i=C_0P_{F_0}w_i
 \qquad(1\le i\le N).
\]
Since $u\ne0$, Lemma~\ref{lem:generic-wu} excludes $\dim F>d_*$.
Together with $W\subset F$ this yields $F=W$.  Moreover $F_0$ is
nonzero.  Otherwise $F_\infty=W$ and $P_W^Nu=0$, contradicting
the injectivity in \eqref{eq:projection-U}.

Extend $C_0$ by zero on $F_\infty$ to obtain a nonzero symmetric
operator $C$ on $W$.  We next verify that $C$ is an inverse limit
covered by Lemma~\ref{lem:inverse-limits}.  The two groups defining
$F$ consist of exactly $d_*$ eigenvalues, bounded away from zero
in absolute value, and the span of their eigenvectors converges to $W$.
All remaining eigenvalues tend to zero.  Let
\[
 A_n=P_WY_n\big|_W
\]
be the compression to $W$.  To identify its inverse, first keep only
the $d_*$ eigenvalues in the two groups.  In a fixed orthonormal
basis of $W$, their contribution to $A_n$ is
\[
 A_n^{\mathrm{large}}=S_n\Lambda_nS_n^T,
\]
where the columns of $S_n$ are the projected eigenvectors and
$\Lambda_n$ is the diagonal matrix of these eigenvalues.
The matrices $S_n$ converge to an orthogonal matrix, while
$\Lambda_n^{-1}$ converges to the diagonal matrix of the
reciprocal finite limits and zeros for the divergent eigenvalues.
Consequently
\[
 (A_n^{\mathrm{large}})^{-1}\longrightarrow C.
\]
The contribution of the remaining eigenvalues tends to zero in
norm.  Since the displayed inverses are bounded, the identity
\[
 A_n^{-1}
 =\bigl(I+(A_n^{\mathrm{large}})^{-1}
                  (A_n-A_n^{\mathrm{large}})\bigr)^{-1}
    (A_n^{\mathrm{large}})^{-1}
\]
shows that $A_n$ is invertible for all sufficiently large $n$ and
$A_n^{-1}\to C$.  Lemma~\ref{lem:inverse-limits} now gives
$[C]\in\mathcal C$.

The projection identities above say $P_Wu_i=Cw_i$.  Since $u\in U$
and $U\subset w^\perp$, we obtain
\[
 0=\langle w,u\rangle
   =\sum_{i=1}^N\langle w_i,P_Wu_i\rangle
   =\sum_{i=1}^N w_i^TCw_i.
\]
This contradicts \eqref{eq:inverse-pairings} and proves the lemma.
\end{proof}

The preceding separation has a useful quadratic consequence.  The
set $\mathcal S_U$ is closed under multiplication by any real
scalar.  Hence its closed set of directions on the unit sphere
contains neither $w$ nor $-w$.  Compactness gives a number
$\delta>0$ such that
\[
 \|P_{w^\perp}z\|\ge\delta\|z\|
 \qquad(z\in\mathcal S_U).
\]
Here the projections act on $H^N$, and $w$ has unit length.  Define
the symmetric operator
\[
 g=\delta^{-2}P_{w^\perp}-P_{\R w}.
\]
For $z\in\mathcal S_U$, its quadratic form satisfies
$z^Tgz\ge\|z\|^2-\|P_{\R w}z\|^2\ge0$, whereas $w^Tgw=-1$.
We have therefore obtained
\begin{equation}
\label{eq:g-properties}
 w^Tgw<0,\qquad
 \bigl((I_N\otimes Y)u\bigr)^T
 g\bigl((I_N\otimes Y)u\bigr)\ge0
 \quad(u\in U,\ Y\text{ Hankel}).
\end{equation}

We can now define the quadratic matrix used in the rest of the
proof.  Use the monomial coordinates to identify
$\R[x_0,\ldots,x_7]_4^*$ with $\R^m$, where $m=330$.
For such a functional $y$, let $Y(y)$ be its Hankel matrix on $H$,
so $Y(y)(p,p')=y(pp')$, and set
\[
 G(y)=I_N\otimes Y(y).
\]
Choose a basis of $U$ and also write $U$ for the $hN$-by-$r$ matrix
whose columns are these basis vectors.  Define
\begin{equation}
\label{eq:constructed-Q}
 Q(y)=U^TG(y)gG(y)U.
\end{equation}
The entries of $G(y)$ are linear in $y$, so $Q$ is a real symmetric
$r$-by-$r$ matrix with homogeneous quadratic entries.  For every
$a\in\R^r$, the vector $Ua$ belongs to the subspace $U$, and
\eqref{eq:g-properties} gives
\[
 a^TQ(y)a=(G(y)Ua)^Tg(G(y)Ua)\ge0.
\]
Thus $Q(y)\succeq0$ for every real $y$.  The data retained for the
next section are $T$, $w$, $U$, and $g$, together with this
admissible map $Q$; the strict inequality at $w$ in
\eqref{eq:g-properties} will produce the obstruction to a finite
semidefinite lift.
